\section{The API Is Not the Air}
\label{sec:ble-api-not-air}

This section could not have been written before the board was flashing, because everything in it is a
measurement that contradicts a success message. It is the part of Bluetooth bring-up that no header
documents, because the header is describing the API and the failure is between the API and the antenna.

\subsection{Four Successes, One Silence}

The first working attempt reported, on every link of the chain, that it was fine:

\begin{minted}{text}
RAIL_Init                 handle returned
RAIL_BLE_Init             void
RAIL_BLE_ConfigPhy1Mbps   status 0
RAIL_ConfigTxPower        status 0
RAIL_SetTxPowerDbm(100)   status 0
RAIL_StartTx              status 0
RAIL_EVENT_TX_PACKET_SENT 39529 times   <- counted in the radio's own interrupt
\end{minted}

and no receiver heard anything. The counter in the last line is the one that makes this worth writing
down: it is not the program counting its own calls, it is the radio's interrupt handler reporting that
the modem completed a packet. The packet was built correctly too, and read back off the chip as

\begin{minted}{text}
00 14 | 00 11 22 33 44 55 | 02 01 06 | 0A 09 "xiao-mg24"
\end{minted}

which is a well-formed \codeinline{ADV\_IND}: public address, general-discoverable flags, complete local
name. A scanner at one hundred percent duty cycle heard four other advertisers in the same room, up to
262 reports each, and this board zero times.

The API was not lying about what it was asked. It was reporting completion of a transmission that had
no power behind it.

\subsection{The Crystal Has To Be Selected, Not Merely Started}

The radio is clocked from the 39 MHz crystal. Starting it is half the job:

\begin{minted}{c}
HFXO0->CTRL |= HFXO_CTRL_FORCEEN;                       /* the oscillator runs */
CMU_ClockSelectSet(cmuClock_SYSCLK, cmuSelect_HFXO);   /* and now it is the clock */
\end{minted}

With only the first line, \codeinline{RAIL\_GetRadioClockFreqHz()} reports 19000 kHz --- the internal
oscillator --- and the BLE channel configuration is refused with
\codeinline{RAIL\_ASSERT\_INVALID\_XTAL\_FREQUENCY}. That assert is not a complaint about the crystal; it
is the library saying that the frequency it was given and the frequency it is running at are different
numbers, and it is right.

Two smaller traps live in the same paragraph. \codeinline{CMU\_OscillatorEnable} in this emlib is
documented as a dummy function kept for backward compatibility: it compiles, it is callable, and it
does nothing. And \codeinline{CMU\_ClockFreqGet(cmuClock\_HFXO)} returns zero because the function's
switch has no case for that clock, so a zero there says nothing at all about the oscillator.

\subsection{The Amplifier Needs Its Curves, and the Call Is One Line}

The gap between the modem and the antenna was a missing initialisation, and the library documents it in
a comment inside its own conversion function: without \codeinline{RAIL\_InitTxPowerCurves} the
dBm-to-raw conversion finds no curve and returns 255, its ``invalid'' value. Silicon Labs' entry point
for that is one call:

\begin{minted}{c}
sl_rail_util_pa_init();                  /* loads the PA curves into RAIL */
RAIL_ConfigTxPower(handle, sl_rail_util_pa_get_tx_power_config_2p4ghz());
RAIL_SetTxPowerDbm(handle, 100);
\end{minted}

The measurement that identifies this failure is a read-back, not a status:

\begin{minted}{text}
                            before                      after sl_rail_util_pa_init()
RAIL_SetTxPowerDbm(100)     0 (success)                 0 (success)
RAIL_GetTxPowerDbm()        -32768  (invalid power)     99   (10 dBm)
RAIL_GetTxPower()           0       (no level)          0x24
\end{minted}

A radio with no valid amplifier level completes every packet it is asked to send. The success in the
first row is true and worthless; the two read-backs are the only rows that describe the air.

\subsection{The Vector Table Had No Radio In It}

Before any of that could be measured, the first transmission faulted: \codeinline{HardFault},
\codeinline{CFSR=0x00000100} (imprecise bus error), at \codeinline{RAIL\_StartTx}. The cause was in the
PAC, not in the radio: it declared 46 of the device's 74 interrupts, and every radio one was among the
missing --- \codeinline{AGC}, \codeinline{BUFC}, \codeinline{FRC\_PRI}, \codeinline{FRC},
\codeinline{MODEM}, \codeinline{PROTIMER}, \codeinline{RAC\_RSM}, \codeinline{RAC\_SEQ},
\codeinline{SYNTH}, \codeinline{RFECA0/1}. Those slots held zeroes in the vector table, so the first
radio interrupt jumped to address zero.

Two further layers hide the same defect, and both were measured:

\begin{itemize}
    \item The vendor names its handlers \codeinline{<NAME>\_IRQHandler}. A vector entry saying
          \codeinline{FRC} falls short of the blob's \codeinline{FRC\_IRQHandler} by the suffix.
    \item A \codeinline{PROVIDE(FRC = DefaultHandler)} \emph{satisfies} the vector's reference, so the
          linker never pulls the archive member that defines the real handler. With the aliases in
          place the program halted inside \codeinline{DefaultHandler}. Linking the RAIL blob with
          \codeinline{+whole-archive} keeps both: a PAC example with no radio still links, and a binary
          with the radio gets RAIL's handlers.
\end{itemize}

\subsection{The Instrument Matters More Than the Hypothesis}

Two instruments were needed and neither was the obvious one.

\codeinline{hcitool lescan} and \codeinline{bluetoothctl} printed nothing on this host's adapter,
including while the host itself was advertising, so both would have been read as ``nothing on air''.
The scan that works is set up through raw HCI with the window equal to the interval --- one hundred
percent duty cycle --- and the reports are read back out of a \codeinline{btmon} capture, which sees
other advertisers the command-line tools never printed.

\codeinline{RAIL\_GetRadioState} returning idle after each transmission proves only that the state
machine moved. The value that identified the fault was the power read-back; the value that proved the
modem was working all along was the send-complete event, counted in an interrupt.

\subsection{What the Working Sequence Is}

\begin{minted}{c}
emlib_hfxo_start(); emlib_clock_select_hfxo();
RAIL_Init(&cfg, NULL);                        /* with an events callback */
RAIL_BLE_Init(handle);
RAIL_BLE_ConfigPhy1MbpsViterbi(handle);       /* only the Viterbi 1M config exists here */
sl_rail_util_pa_init();                       /* the PA curves: without this, silence */
RAIL_ConfigTxPower(handle, sl_rail_util_pa_get_tx_power_config_2p4ghz());
RAIL_SetTxPowerDbm(handle, 100);
RAIL_ConfigData(&packetMode); RAIL_SetTxFifo(txBuf, 0, 64);
for (;;) {
  channel = 37 + (n % 3);                     /* the three advertising channels */
  RAIL_BLE_ConfigChannelRadioParams(handle, 0x555555, 0x8E89BED6, channel, false);
  RAIL_PrepareChannel(handle, channel);
  RAIL_WriteTxFifo(pdu, pduLen, true);        /* the full PDU, header included */
  RAIL_StartTx(handle, channel, RAIL_TX_OPTIONS_DEFAULT, NULL);
}
\end{minted}

The radio parameters are re-applied per channel because the whitening seed is the channel number.
\end{document} is not needed; the loop is the whole application.

With that sequence the board advertises, and the proof is a capture from another machine:

\begin{minted}{text}
Address: 55:44:33:22:11:00 (OUI 55-44-33)
Flags: 0x06   LE General Discoverable Mode, BR/EDR Not Supported
Name (complete): xiao-mg24
RSSI: -83 dBm
\end{minted}

\subsection{The Rule This Chapter Adds to the Book}

A radio that reports success is not a radio that radiates. Every link in this chain --- the channel
configuration, the transmission, the modem, the interrupt count --- can be green while the antenna is
silent, and the only measurement that distinguishes them is one taken \emph{outside the chip}: a
receiver, and a read-back of the amplifier's own idea of its power. The build was correct for a long
time before it was true.
